% !TEX root = ../../../main.tex
\section{Mean value theorems}
\label{sec:analysis-i:mean-value-theorems}

% Source: images/LEC06/IMG_9354--IMG_9356.
% Rephotographed as images/LEC06/IMG_9415--IMG_9417.
\begin{theorem}[Rolle's theorem]
  \label{thm:analysis-i:rolle}
  Let $f$ be continuous on $[a,b]$ and differentiable on $(a,b)$,
  where $a<b$. If $f(a)=f(b)$, then there exists
  $c\in(a,b)$ such that $f'(c)=0$.
\end{theorem}
\begin{proof}
  If $f$ is constant, every interior point has derivative zero.
  Otherwise, by \autoref{thm:analysis-i:extreme-value}, $f$ attains a
  maximum and a minimum. Since the endpoint values agree and $f$ is
  nonconstant, at least one of these extrema is attained at an interior
  point $c$. Then \autoref{thm:analysis-i:interior-extremum} gives
  $f'(c)=0$.
\end{proof}
\historynote*{Rolle's 1691 \emph{D\'emonstration} concerns solving
polynomial equations~\cite{Rolle1691}. The formulation here is a modern
statement for differentiable functions, not a verbatim seventeenth-century
theorem. Equal endpoint heights force a stationary point, making this a
useful bridge from extrema to mean-value arguments.}

The following polynomial question can be approached
by constructing a polynomial whose derivative is $P$ and comparing its
endpoint values. No integration theorem is required.
% !TEX root = ../../hw07.tex
% Source: MA-GY-6213-hw7 (1).pdf, page 1, Q3.
\begin{exercise}
  \label{ex:hw07:polynomial-rolle}
  Let $P(x)=\sum_{k=0}^n a_kx^k$ satisfy
  \[
    a_0+\frac{a_1}{2}+\cdots+\frac{a_n}{n+1}=0.
  \]
  Show that there is a point $x\in(0,1)$ such that $P(x)=0$.
  \solutionlink{hw07:polynomial-rolle}
\end{exercise}


\begin{theorem}[Mean value theorem]
  \label{thm:analysis-i:mean-value}
  Let $a<b$. Suppose $f$ is continuous on $[a,b]$
  and differentiable on $(a,b)$. There exists $c\in(a,b)$ such that
  \[
    f'(c)=\frac{f(b)-f(a)}{b-a}.
  \]
\end{theorem}
\historynote*{Cauchy's 1823 \emph{R\'esum\'e des le\c cons} develops
differential calculus through derivatives and finite
increments~\cite{Cauchy1823}. The mean-value method relates a finite
change to a local slope, turning derivative information into bounds and
monotonicity. This citation documents a historical treatment, without
claiming sole discovery of the modern formulation.}

\begin{proof}
  Subtract the secant line by defining
  \[
    H(t)=f(t)-f(a)-\frac{f(b)-f(a)}{b-a}(t-a).
  \]
  This function is continuous on $[a,b]$, differentiable on $(a,b)$,
  and satisfies $H(a)=H(b)=0$. By \autoref{thm:analysis-i:rolle}, some
  $c\in(a,b)$ satisfies $H'(c)=0$. Therefore
  \[
    0=f'(c)-\frac{f(b)-f(a)}{b-a},
  \]
  which is the required equality.
\end{proof}
The geometry of both results is shown in
\autoref{fig:analysis-i:mean-value-geometry}.
\begin{figure*}
  \centering
  % Source sketches: LEC06/IMG_9354, Rolle and mean value theorem.
\begin{tikzpicture}[interval diagram,x=1cm,y=1cm]
  \begin{scope}
    \draw[->] (0,0)--(4.4,0) node[right=4pt] {$x$};
    \draw[->] (0,0)--(0,3) node[above=4pt] {$y$};
    \draw[domain=.6:3.8,samples=50]
      plot (\x,{2.45-.45*(\x-2.2)^2});
    \draw[gray,dashed] (.6,0)--(.6,1.298)--(3.8,1.298)--(3.8,0);
    \draw[gray,dashed] (2.2,0)--(2.2,2.45);
    \draw[gray] (1.55,2.45)--(2.85,2.45);
    \fill (.6,1.298) circle[radius=1.6pt];
    \fill (3.8,1.298) circle[radius=1.6pt];
    \fill (2.2,2.45) circle[radius=1.6pt];
    \foreach \x/\lab in {.6/a,2.2/c,3.8/b}
      \node[below=5pt] at (\x,0) {$\lab$};
    \node[above=9pt] at (2.2,3) {Rolle's theorem};
  \end{scope}
  \begin{scope}[xshift=6cm]
    \draw[->] (0,0)--(4.4,0) node[right=4pt] {$x$};
    \draw[->] (0,0)--(0,3) node[above=4pt] {$y$};
    \draw[domain=.6:3.8,samples=50] plot (\x,{.6+.14*\x*\x});
    \draw[gray] (.6,.6504)--(3.8,2.6216);
    \draw[gray,dashed] (.6,0)--(.6,.6504);
    \draw[gray,dashed] (3.8,0)--(3.8,2.6216);
    \draw[gray,dashed] (2.2,0)--(2.2,1.2776);
    \draw (1.35,.754)--(3.05,1.8012);
    \fill (.6,.6504) circle[radius=1.6pt];
    \fill (3.8,2.6216) circle[radius=1.6pt];
    \fill (2.2,1.2776) circle[radius=1.6pt];
    \foreach \x/\lab in {.6/a,2.2/c,3.8/b}
      \node[below=5pt] at (\x,0) {$\lab$};
    \node[above=9pt] at (2.2,3) {Mean value theorem};
  \end{scope}
\end{tikzpicture}

  \caption[Rolle and mean-value geometry.]{On the
    left, equal endpoint heights give a horizontal tangent. On the right,
    the tangent at $c$ is parallel to the secant through the endpoints.
    These are schematic examples of the two theorems.}
  \label{fig:analysis-i:mean-value-geometry}
\end{figure*}

\subsection{An average-slope interpretation}
% IMG_9354's integration picture motivates MVT; it is not its proof.
If $f'$ is continuous on $[a,b]$, the fundamental theorem of calculus,
assumed here and reserved for
\autoref{sec:analysis-i:fundamental-theorem-of-calculus}, gives
\[
  \frac1{b-a}\int_a^b f'(x)\,dx=\frac{f(b)-f(a)}{b-a}.
\]
Thus the mean value theorem says that this average slope is attained
at an interior point. The continuity of $f'$ is an additional hypothesis
for this integral interpretation. The mean value theorem itself only
requires differentiability on the open interval. A schematic illustration
of average height appears in
\autoref{fig:analysis-i:average-height}.
\begin{marginfigure}
  \centering
  % Schematic average-height sketch from LEC06/IMG_9354.
\begin{tikzpicture}[interval diagram,x=1cm,y=1cm]
  \draw[->] (0,0)--(4.1,0) node[right=4pt] {$x$};
  \draw[->] (0,0)--(0,2.55) node[above=4pt] {$\phi$};
  \draw[domain=.45:3.55,samples=100]
    plot (\x,{1.55+.32*cos(360*(\x-.45)/1.55)});
  \draw[gray,dashed] (.45,1.55)--(3.55,1.55);
  \draw[gray,dashed] (.45,0)--(.45,1.87);
  \draw[gray,dashed] (3.55,0)--(3.55,1.87);
  \node[below=5pt] at (.45,0) {$a$};
  \node[below=5pt] at (3.55,0) {$b$};
  \node[above=5pt] at (2,2.05) {$\overline\phi$};
\end{tikzpicture}

  \caption[Average height of a continuous graph.]{A
    schematic continuous function $\phi$ and its average height
    $\overline\phi=(b-a)^{-1}\int_a^b\phi(x)\,dx$ (dashed).}
  \label{fig:analysis-i:average-height}
\end{marginfigure}

\subsection{Derivative signs and monotonicity}
\begin{corollary}
  \label{thm:analysis-i:derivative-monotonicity}
  Let $f$ be differentiable on an open interval $I$.
  If $f'(x)\ge0$ for every $x\in I$, then $f$ is nondecreasing on $I$.
  If $f'(x)>0$ everywhere, it is strictly increasing.
  The analogous statements with reversed inequalities also hold.
\end{corollary}
\begin{proof}
  For $x_1<x_2$ in $I$, differentiability implies continuity on
  $[x_1,x_2]$. \autoref{thm:analysis-i:mean-value} gives a $c\in(x_1,x_2)$ with
  $f(x_2)-f(x_1)=f'(c)(x_2-x_1)$. This difference is nonnegative when
  $f'\ge0$, and positive when $f'>0$. Apply the same argument to $-f$
  for the reversed inequalities.
\end{proof}
% IMG_9354's nondecreasing shorthand mixes strict and non-strict signs.

The next exercise links the mean value theorem to Lipschitz continuity
from \autoref{sec:analysis-i:uniform-continuity}. Its iteration also uses
geometric-series tail bounds and the completeness of $\R$ from
\autoref{sec:analysis-i:cauchy-sequences}.
% !TEX root = ../../hw07.tex
% Source: MA-GY-6213-hw7 (1).pdf, page 1, Q1, (i)--(iii).
\begingroup
\renewcommand{\labelenumi}{(\roman{enumi})}
\begin{exercise}
  \label{ex:hw07:contractive-fixed-point}
  Let $f$ be a differentiable function over $\R$.
  \begin{enumerate}
    \item Show that if $f'(x)$ is a bounded function, then $f$ is a
      Lipschitz function which satisfies
      \[
        |f(x)-f(y)|\le L|x-y|,\qquad x,y\in\R,
      \]
      where $L>0$ is a constant called the Lipschitz constant.
    \item If the Lipschitz constant in (i) satisfies $L<1$, the function
      is called ``contractive''. Prove that a contractive function must
      have a fixed point $x^*$ such that $x^*=f(x^*)$.
      (Hint: Take any $x_0\in\R$ and consider the iterative sequence
      $x_n=f(x_{n-1})$ for $n=1,2,\ldots$.)
    \item Show that the fixed point $x^*$ found in (ii) is unique.
  \end{enumerate}
  \solutionlink{hw07:contractive-fixed-point}
\end{exercise}
\endgroup


Derivative signs can also identify a sharp global
bound on a half-line. The next question asks for both the bound and
the equality case.
% !TEX root = ../../hw07.tex
% Source: MA-GY-6213-hw7 (1).pdf, page 2, Q5.
\begin{exercise}
  \label{ex:hw07:rational-maximum}
  Let $a>1$. Show that for all $x\ge0$, we have
  \[
    \frac1{1+x}-\frac1{1+ax}\le\frac{\sqrt a-1}{\sqrt a+1},
  \]
  with equality only for $x=1/\sqrt a$.
  \solutionlink{hw07:rational-maximum}
\end{exercise}


\subsection{Why logarithms preserve monotonicity}

For a positive function, taking a logarithm changes the heights of its
graph while preserving the order of those heights. This is why its
logarithm can serve as a proxy for studying increase and decrease.
The essential property is that $\ln$ is strictly increasing on
$(0,\infty)$. The exponential and logarithm used here are constructed
in \autoref{thm:analysis-i:exponential-properties} and
\autoref{thm:analysis-i:logarithm-real-powers}.

\begin{figure*}[htbp]
  \centering
  \input{\analysisroot/figures/interval-styles.tex}
\begin{tikzpicture}[interval diagram,x=1.2cm,y=.9cm]
  % All panels have identical units and the same displayed window.
  \foreach \shift/\heading in
    {0/Increasing,3.75/Decreasing,7.5/A peak,11.25/A valley} {
    \begin{scope}[shift={(\shift cm,0)}]
      \node at (0,3.9) {\heading};
      \draw[->] (-1.15,0)--(1.25,0) node[right=4pt] {$x$};
      \draw[->] (0,-1.5)--(0,3.2) node[above=4pt] {$y$};
      \draw[gray,dotted] (-1.12,1)--(1.1,1);
      \node[left=4pt] at (-1.12,1) {$1$};
      \draw (-1,0)--(-1,-.06) node[below=5pt] {$-1$};
      \draw (1,0)--(1,-.06) node[below=5pt] {$1$};
    \end{scope}
  }
  \begin{scope}
    \draw[<->,domain=-1.05:1.05,samples=70] plot (\x,{exp(\x)});
    \draw[gray,dashed,<->] (-1.05,-1.05)--(1.05,1.05);
    \node at (0,-2.05) {$f(x)=e^x$};
    \node at (0,-2.65) {$\ln f(x)=x$};
  \end{scope}
  \begin{scope}[shift={(3.75cm,0)}]
    \draw[<->,domain=-1.05:1.05,samples=70] plot (\x,{exp(-\x)});
    \draw[gray,dashed,<->] (-1.05,1.05)--(1.05,-1.05);
    \node at (0,-2.05) {$f(x)=e^{-x}$};
    \node at (0,-2.65) {$\ln f(x)=-x$};
  \end{scope}
  \begin{scope}[shift={(7.5cm,0)}]
    \draw[<->,domain=-1.05:1.05,samples=70] plot (\x,{exp(-\x*\x)});
    \draw[gray,dashed,<->,domain=-1.05:1.05,samples=70]
      plot (\x,{-\x*\x});
    \node at (0,-2.05) {$f(x)=e^{-x^2}$};
    \node at (0,-2.65) {$\ln f(x)=-x^2$};
  \end{scope}
  \begin{scope}[shift={(11.25cm,0)}]
    \draw[<->,domain=-1.05:1.05,samples=70] plot (\x,{1+\x*\x});
    \draw[gray,dashed,<->,domain=-1.05:1.05,samples=70]
      plot (\x,{ln(1+\x*\x)});
    \node at (0,-2.05) {$f(x)=1+x^2$};
    \node at (0,-2.65) {$\ln f(x)=\ln(1+x^2)$};
  \end{scope}
  \draw (3.6,4.65)--(4.05,4.65);
  \node[anchor=west] at (4.18,4.65) {$f$};
  \draw[gray,dashed] (5.35,4.65)--(5.8,4.65);
  \node[anchor=west] at (5.93,4.65) {$\ln f$};
\end{tikzpicture}

  \caption[Logarithms preserve order.]{Solid curves
    are $f$ and dashed curves are $\ln f$, on identical coordinate
    scales. Dotted guides mark $y=1$. The logarithm preserves increase,
    decrease, and the inputs at peaks and valleys. Values and curvature
    can change. Only a window of each graph is shown.}
  \label{fig:analysis-i:logarithm-monotonicity}
\end{figure*}

\Needspace{13\baselineskip}
\begin{proposition}[Logarithms preserve order and monotonicity]
  \label{thm:analysis-i:logarithmic-monotonicity}
  Let $f:I\to(0,\infty)$, where $I$ is an interval, and put
  $H(x)=\ln(f(x))$. For any $s,t\in I$,
  \[
    f(s)<f(t)\quad\Longleftrightarrow\quad H(s)<H(t).
  \]
  The same equivalence holds with $\le$ or $=$. Consequently, $f$
  and $H$ are increasing or decreasing on exactly the same intervals,
  with the same strict or non-strict sense. They also attain local
  and global maxima and minima at the same inputs.
\end{proposition}
\begin{proof}
  First take $0<a<b$. The logarithm is continuous on $[a,b]$ and
  differentiable on $(a,b)$. The mean value theorem gives some
  $\xi\in(a,b)$ with
  \[
    \ln b-\ln a=\frac{b-a}{\xi}>0.
  \]
  Thus $\ln$ preserves strict inequalities between positive values.
  Conversely, if $H(s)<H(t)$, strict increase of the exponential gives
  $f(s)=e^{H(s)}<e^{H(t)}=f(t)$. Its inverse identities also preserve
  equality, so the non-strict comparison follows.

  Applying these comparisons to every pair $s<t$ proves the
  monotonicity statements. Applying them to the comparison of one
  value with all nearby values, or all values in $I$, proves the
  statements about extrema.
\end{proof}

When $f$ is differentiable, the chain rule gives a second way to see
the connection,
\[
  H'(x)=\frac{f'(x)}{f(x)}.
\]
Since $f(x)>0$, the derivatives $H'(x)$ and $f'(x)$ have the same sign
and vanish at the same inputs. A simpler formula for $H'$ can therefore
settle the sign of $f'$ without differentiating a complicated power
directly. For example, if $f(x)=g(x)^{h(x)}$ with $g(x)>0$, then
\[
  H(x)=h(x)\ln g(x).
\]
The exponent has become a factor.

\autoref{fig:analysis-i:logarithm-monotonicity} compares four positive
functions with their logarithms using matching axes and identical scales.
Each point $(x,f(x))$ becomes $(x,\ln(f(x)))$. The input stays fixed
while the vertical coordinate changes. The slope $1/y$ of the logarithm
compresses differences between large positive outputs and magnifies
differences between outputs close to zero.
The increasing and decreasing examples retain their direction.
The other two retain their peak and valley at $x=0$.
The heights change, and curvature need not be preserved.
In particular, $f(x)=1$ becomes $H(x)=0$, while $0<f(x)<1$ gives
$H(x)<0$. A negative logarithm is allowed. It is the argument of
the logarithm that must be positive throughout the interval.

\begin{example}[Increasing exponential approximants]
  \label{ex:analysis-i:increasing-euler-approximants}
  Consider $F(x)=(1+1/x)^x$ for $x\ge1$. We will
  show that $F$ is strictly increasing. Since $F>0$,
  \autoref{thm:analysis-i:logarithmic-monotonicity} lets us study its
  logarithm with exactly the same monotonicity. Taking the logarithm
  also turns the varying exponent into a product. Put
  \[
    H(x)=\ln F(x)=x\ln(1+1/x).
  \]
  The left panel of \autoref{fig:analysis-i:increasing-euler-approximants}
  shows the positive gap that will control the derivative. The right
  panel illustrates the resulting increase of $F$.
  \begin{figure*}[htbp]
    \centering
    \input{\analysisroot/figures/interval-styles.tex}
\begin{tikzpicture}[interval diagram]
  \begin{scope}[x=4.3cm,y=5.2cm]
    \node[anchor=west] at (0,.98) {The gap controls the derivative};
    \draw[->] (0,0)--(1.19,0) node[right=5pt] {$u$};
    \draw[->] (0,0)--(0,.78) node[above=5pt] {$y$};
    \node[below left=5pt] at (0,0) {$0$};
    \fill[black!6]
      plot[domain=0:1,samples=70] (\x,{ln(1+\x)})
      -- plot[domain=1:0,samples=70] (\x,{\x/(1+\x)}) -- cycle;
    \draw[domain=0:1,samples=70] plot (\x,{ln(1+\x)});
    \draw[gray,dashed,domain=0:1,samples=70]
      plot (\x,{\x/(1+\x)});
    \fill (1,{ln(2)}) circle[radius=1.6pt];
    \fill[gray] (1,.5) circle[radius=1.6pt];
    \draw (1,0)--(1,-.015) node[below=5pt] {$1$};
    \draw[gray,dotted] (.8,0)--(.8,{.8/1.8});
    \node[below=5pt] at (.8,0) {$u_0$};
    \draw[<->] (.8,{.8/1.8})--(.8,{ln(1.8)});
    \draw[gray] (.83,.516)--(1.10,.516);
    \node[right=5pt] at (1.10,.516) {$\psi(u_0)$};
    \draw (0,-.17)--(.14,-.17);
    \node[anchor=west] at (.18,-.17) {$y=\ln(1+u)$};
    \draw[gray,dashed] (0,-.29)--(.14,-.29);
    \node[anchor=west] at (.18,-.29) {$y=u/(1+u)$};
  \end{scope}
  \begin{scope}[shift={(7.4cm,0)},x=.56cm,y=1.35cm]
    \node[anchor=west] at (0,3.775) {The approximants increase};
    \draw[->] (0,0)--(8.8,0) node[right=5pt] {$x$};
    \draw[->] (0,0)--(0,3) node[above=5pt] {$y$};
    \node[below left=5pt] at (0,0) {$0$};
    \draw[gray,dashed] (0,{exp(1)})--(8.45,{exp(1)});
    \node[left=5pt] at (0,{exp(1)}) {$e$};
    \draw (0,2)--(-.12,2) node[left=5pt] {$2$};
    \draw[->,domain=1:8.45,samples=80]
      plot (\x,{exp(\x*ln(1+1/\x))});
    \foreach \n in {1,2,4,8} {
      \fill (\n,{exp(\n*ln(1+1/\n))}) circle[radius=1.6pt];
      \draw (\n,0)--(\n,-.06) node[below=5pt] {$\n$};
    }
    \draw (0,-.655)--(1.08,-.655);
    \node[anchor=west] at (1.38,-.655) {$F(x)=(1+1/x)^x$};
    \draw[gray,dashed] (0,-1.12)--(1.08,-1.12);
    \node[anchor=west] at (1.38,-1.12) {limiting level $e$};
  \end{scope}
\end{tikzpicture}

    \caption[The gap behind increasing approximants.]{%
      The shaded gap on the left is
      $\psi(u)=\ln(1+u)-u/(1+u)$, with $u_0=4/5$ marked.
      Its positivity gives $H'(x)=\psi(1/x)>0$.
      The right panel shows the actual graph of $F$, with selected
      integer inputs marked. The dashed level is $e$, whose role as
      the limit is proved in
      \autoref{ex:analysis-i:euler-approximant-limit}.}
    \label{fig:analysis-i:increasing-euler-approximants}
  \end{figure*}
  \begin{proof}
  The product and chain rules give the first line below. The second
  uses $x/[x^2(1+1/x)]=1/(x+1)$.
  \[
    \begin{aligned}
      H'(x)&=\ln(1+1/x)
        +x\frac{1}{1+1/x}\left(-\frac1{x^2}\right)\\
        &=\ln(1+1/x)-\frac1{x+1}.
    \end{aligned}
  \]
  % IMG_9355 stops here with "Finish off later".
  The sign depends on which of the two positive terms is larger.
  Write $u=1/x$, so $1/(x+1)=u/(1+u)$, and define
  \[
    \psi(u)=\ln(1+u)-\frac{u}{1+u},\qquad u\ge0.
  \]
  Thus $H'(x)=\psi(1/x)$. Geometrically, $\psi(u)$ is the vertical
  gap between the two curves in the left panel.

  The gap is zero at $u=0$. For $u>0$, its derivative is
  \[
    \psi'(u)=\frac1{1+u}-\frac1{(1+u)^2}
      =\frac{u}{(1+u)^2}>0.
  \]
  To turn this derivative information into a strict inequality, fix
  $u>0$. The function $\psi$ is continuous on $[0,u]$ and differentiable
  on $(0,u)$. By \autoref{thm:analysis-i:mean-value}, some
  $\xi\in(0,u)$ satisfies
  \[
    \psi(u)-\psi(0)=u\psi'(\xi)
      =\frac{u\xi}{(1+\xi)^2}>0.
  \]
  Thus $\psi(u)>0$ for $u>0$. Taking $u=1/x$ shows that
  $H'(x)=\psi(1/x)>0$ for $x>1$.

  Finally, take $1\le s<t$. The function $H$ is continuous on $[s,t]$
  and differentiable on $(s,t)$. The mean value theorem gives
  $H(t)-H(s)=H'(c)(t-s)>0$ for some $c\in(s,t)$.
  Since the exponential is strictly increasing, the following comparison
  proves that $F$ is strictly increasing on $[1,\infty)$.
  \[
    F(s)=e^{H(s)}<e^{H(t)}=F(t). \qedhere
  \]
  \end{proof}
\end{example}

The next question changes the exponent by one.
The corresponding limit can be studied after
\autoref{ex:analysis-i:euler-approximant-limit}, where the logarithmic
quotient is evaluated. Its monotonicity can already be studied here.
% !TEX root = ../../hw07.tex
% Source: MA-GY-6213-hw7 (1).pdf, page 2, Q6. Domain unspecified in PDF.
% Source spelling "descreases" regularised, without a mathematical change.
\begin{exercise}
  \label{ex:hw07:decreasing-euler-approximants}
  Show that the function
  \[
    \left(1+\frac1x\right)^{x+1}
  \]
  decreases and find the limit
  \[
    \lim_{x\to\infty}\left(1+\frac1x\right)^{x+1}.
  \]
  \marginnote{The source does not specify a domain.
    Take $x>0$ for this positive-base real-power problem.}
  \solutionlink{hw07:decreasing-euler-approximants}
\end{exercise}


\subsection{The extended mean value theorem}
\begin{theorem}[Cauchy's mean value theorem]
  \label{thm:analysis-i:cauchy-mean-value}
  Let $a<b$, and suppose $f$ and $g$ are continuous
  on $[a,b]$ and differentiable on $(a,b)$. There exists $c\in(a,b)$
  such that
  \[
    f'(c)\bigl(g(b)-g(a)\bigr)
      =g'(c)\bigl(f(b)-f(a)\bigr).
  \]
\end{theorem}
Applying the ordinary mean value theorem separately
to $f$ and $g$ gives two possibly different points. It does not show that
they coincide. The extended theorem asserts a single common point.
Its proof uses one auxiliary function.
% IMG_9355 incorrectly identifies the two separate MVT points. IMG_9356
% leaves the extended theorem's proof unfinished. The proof below is an
% editorial addition requested by the author.
\begin{proof}
  Define
  \[
    H(t)=\bigl(g(b)-g(a)\bigr)\bigl(f(t)-f(a)\bigr)
      -\bigl(f(b)-f(a)\bigr)\bigl(g(t)-g(a)\bigr).
  \]
  By the hypotheses on $f$ and $g$, the function $H$ is continuous on
  $[a,b]$ and differentiable on $(a,b)$. Also, $H(a)=H(b)=0$.
  Therefore \autoref{thm:analysis-i:rolle} supplies a point $c\in(a,b)$
  with $H'(c)=0$. Differentiating gives
  \[
    \bigl(g(b)-g(a)\bigr)f'(c)
      -\bigl(f(b)-f(a)\bigr)g'(c)=0,
  \]
  which is the required identity. No division by an endpoint difference
  is needed, so the argument also covers zero endpoint differences.
\end{proof}
